\documentclass[10pt,a4paper]{amsart}
\usepackage{amsmath,amssymb,amsthm,graphicx}
\usepackage{verbatim}
\newenvironment{reference}{\comment}{\endcomment}
%\newenvironment{reference}{}{}
\newenvironment{slogan}{\comment}{\endcomment}
\newenvironment{history}{\comment}{\endcomment}

% For commutative diagrams we use Xy-pic
\usepackage[all]{xy}

% We use 2cell for 2-commutative diagrams.
\xyoption{2cell}
% Theorem environments.
%
\theoremstyle{plain}
\newtheorem{theorem}[subsection]{Theorem}
\newtheorem{proposition}[subsection]{Proposition}
\newtheorem{lemma}[subsection]{Lemma}

\theoremstyle{definition}
\newtheorem{definition}[subsection]{Definition}
\newtheorem{example}[subsection]{Example}
\newtheorem{exercise}[subsection]{Exercise}
\newtheorem{situation}[subsection]{Situation}

\theoremstyle{remark}
\newtheorem{remark}[subsection]{Remark}
\newtheorem{remarks}[subsection]{Remarks}

\numberwithin{equation}{subsection}

% Macros
%
\def\lim{\mathop{\mathrm{lim}}\nolimits}
\def\colim{\mathop{\mathrm{colim}}\nolimits}
\def\Spec{\mathop{\mathrm{Spec}}}
\def\Hom{\mathop{\mathrm{Hom}}\nolimits}
\def\Ext{\mathop{\mathrm{Ext}}\nolimits}
\def\SheafHom{\mathop{\mathcal{H}\!\mathit{om}}\nolimits}
\def\SheafExt{\mathop{\mathcal{E}\!\mathit{xt}}\nolimits}
\def\Sch{\mathit{Sch}}
\def\Mor{\mathop{\mathrm{Mor}}\nolimits}
\def\Ob{\mathop{\mathrm{Ob}}\nolimits}
\def\Sh{\mathop{\mathit{Sh}}\nolimits}
\def\NL{\mathop{N\!L}\nolimits}
\def\CH{\mathop{\mathrm{CH}}\nolimits}
\def\proetale{{pro\text{-}\acute{e}tale}}
\def\etale{{\acute{e}tale}}
\def\QCoh{\mathit{QCoh}}
\def\Ker{\mathop{\mathrm{Ker}}}
\def\Im{\mathop{\mathrm{Im}}}
\def\Coker{\mathop{\mathrm{Coker}}}
\def\Coim{\mathop{\mathrm{Coim}}}

% Boxtimes
%
\DeclareMathSymbol{\boxtimes}{\mathbin}{AMSa}{"02}

%
% Macros for moduli stacks/spaces
%
\def\QCohstack{\mathcal{QC}\!\mathit{oh}}
\def\Cohstack{\mathcal{C}\!\mathit{oh}}
\def\Spacesstack{\mathcal{S}\!\mathit{paces}}
\def\Quotfunctor{\mathrm{Quot}}
\def\Hilbfunctor{\mathrm{Hilb}}
\def\Curvesstack{\mathcal{C}\!\mathit{urves}}
\def\Polarizedstack{\mathcal{P}\!\mathit{olarized}}
\def\Complexesstack{\mathcal{C}\!\mathit{omplexes}}
% \Pic is the operator that assigns to X its picard group, usage \Pic(X)
% \Picardstack_{X/B} denotes the Picard stack of X over B
% \Picardfunctor_{X/B} denotes the Picard functor of X over B
\def\Pic{\mathop{\mathrm{Pic}}\nolimits}
\def\Picardstack{\mathcal{P}\!\mathit{ic}}
\def\Picardfunctor{\mathrm{Pic}}
\def\Deformationcategory{\mathcal{D}\!\mathit{ef}}
\usepackage{cleveref}
\crefname{originalreference}{original reference}{original references}
\Crefname{originalreference}{Original reference}{Original references}
\makeatletter
\newlabel{divisors-lemma-pullback-effective-Cartier-defined}{{02OO}{source}{}{sourcepointer-1}{}}
\newlabel{divisors-lemma-pullback-effective-Cartier-defined@cref}{{[originalreference][0][]02OO}{[1][0][]source}{}{}{}}
\makeatother
\crefname{originaldefinition}{original definition}{original definitions}
\Crefname{originaldefinition}{Original definition}{Original definitions}
\makeatletter
\newlabel{definition-strict-normal-crossings}{{0BI9}{source}{}{sourcepointer-2}{}}
\newlabel{definition-strict-normal-crossings@cref}{{[originaldefinition][0][]0BI9}{[1][0][]source}{}{}{}}
\makeatother
\crefname{originalreference}{original reference}{original references}
\Crefname{originalreference}{Original reference}{Original references}
\makeatletter
\newlabel{algebra-lemma-characterize-smooth-over-field}{{00TT}{source}{}{sourcepointer-3}{}}
\newlabel{algebra-lemma-characterize-smooth-over-field@cref}{{[originalreference][0][]00TT}{[1][0][]source}{}{}{}}
\makeatother
\crefname{originalreference}{original reference}{original references}
\Crefname{originalreference}{Original reference}{Original references}
\makeatletter
\newlabel{algebra-lemma-dimension-base-fibre-equals-total}{{00ON}{source}{}{sourcepointer-4}{}}
\newlabel{algebra-lemma-dimension-base-fibre-equals-total@cref}{{[originalreference][0][]00ON}{[1][0][]source}{}{}{}}
\makeatother
\title{Pullback of a strict normal crossings divisor by a smooth morphism is a strict normal crossings divisor}
\author{The Stacks Project Authors}
\date{stacks-project:etale.tex:local-snapshot-20261012; GNU FDL 1.2 or later}
\begin{document}
\maketitle
\noindent\textit{Editorial note.} All delivered mathematical text is the author's own, copied verbatim from the pinned Stacks Project chapter file etale.tex; no statement or proof step was rewritten, reordered or removed.\par\medskip
\noindent\textit{Editorial note.} Span recomputation by label pairing (LANE\_BRIEF 10.2): in this chapter the author's proof environment follows the statement environment as a standalone block; the spec covers the complete environment pair (statement span, proof span) as recorded in the pool candidate.\par\medskip
\ifdefined\hypertarget\hypertarget{sourcepointer-1}{}\fi
\noindent\textit{Source reference.} Original reference 02OO: sources/stacks-chow/tags.txt, source lines 3523--3523. Cross-chapter standard prerequisite cited by the author inside the counted span; it is used as a pointer and its own proof is not part of this item.\par\medskip
\ifdefined\hypertarget\hypertarget{sourcepointer-2}{}\fi
\noindent\textit{Source reference.} Original definition 0BI9: sources/stacks-chow/tags.txt, source lines 14482--14482. Cross-chapter standard prerequisite cited by the author inside the counted span; it is used as a pointer and its own proof is not part of this item.\par\medskip
\ifdefined\hypertarget\hypertarget{sourcepointer-3}{}\fi
\noindent\textit{Source reference.} Original reference 00TT: sources/stacks-chow/tags.txt, source lines 1087--1087. Cross-chapter standard prerequisite cited by the author inside the counted span; it is used as a pointer and its own proof is not part of this item.\par\medskip
\ifdefined\hypertarget\hypertarget{sourcepointer-4}{}\fi
\noindent\textit{Source reference.} Original reference 00ON: sources/stacks-chow/tags.txt, source lines 901--901. Cross-chapter standard prerequisite cited by the author inside the counted span; it is used as a pointer and its own proof is not part of this item.\par\medskip
\begin{lemma}
\label{lemma-smooth-pullback-strict-normal-crossings}
\begin{slogan}
Pullback of a strict normal crossings divisor by a smooth
morphism is a strict normal crossings divisor.
\end{slogan}
Let $X$ be a locally Noetherian scheme. Let $D \subset X$ be a
strict normal crossings divisor. If $f : Y \to X$ is a smooth
morphism of schemes, then the pullback $f^*D$ is a
strict normal crossings divisor on $Y$.
\end{lemma}

\begin{proof}
As $f$ is flat the pullback is defined by
Divisors, Lemma \ref{divisors-lemma-pullback-effective-Cartier-defined}
hence the statement makes sense.
Let $q \in f^*D$ map to $p \in D$. Choose a regular system
of parameters $x_1, \ldots, x_d \in \mathfrak m_p$
and $1 \leq r \leq d$ as in
Definition \ref{definition-strict-normal-crossings}.
Since $f$ is smooth the local ring homomorphism
$\mathcal{O}_{X, p} \to \mathcal{O}_{Y, q}$ is flat
and the fibre ring
$$
\mathcal{O}_{Y, q}/\mathfrak m_p \mathcal{O}_{Y, q} =
\mathcal{O}_{Y_p, q}
$$
is a regular local ring (see for example
Algebra, Lemma \ref{algebra-lemma-characterize-smooth-over-field}).
Pick $y_1, \ldots, y_n \in \mathfrak m_q$ which map to a regular
system of parameters in $\mathcal{O}_{Y_p, q}$.
Then $x_1, \ldots, x_d, y_1, \ldots, y_n$ generate the
maximal ideal $\mathfrak m_q$. Hence $\mathcal{O}_{Y, q}$
is a regular local ring of dimension
$d + n$ by Algebra, Lemma \ref{algebra-lemma-dimension-base-fibre-equals-total}
and $x_1, \ldots, x_d, y_1, \ldots, y_n$
is a regular system of parameters. Since $f^*D$ is cut
out by $x_1 \ldots x_r$ in $\mathcal{O}_{Y, q}$ we conclude
that the lemma is true.
\end{proof}

\end{document}
